% GENERATED FILE. Edit canonical lessons, then run npm run generate:books.
% canonical-source-hash: fnv1a64:18ab5697fc68aabd
% canonical-lessons: MR-W40-i-independent, MR-C40-ithe, MR-R40-asking-column

\chapter{Here, and the Asking Column}
\label{ch:mr-here}
\hlchaptermodality{40}

\begin{chapteropening}
\textbf{By the end of this chapter:} \emph{I can write the independent \mr{इ}, say where I am with ithe, and turn the lines of my own A1 message into the questions they answer.}

Complete MR-R40-asking-column and interrogate your own named-reader message: turn its name, wellbeing and place lines into questions by changing one word each.
\end{chapteropening}

\section[i]{\mr{इ} — short \emph{i} with nothing to lean on}

\label{lesson:MR-W40-i-independent}



\begin{quote}
You hang a short \emph{i} in front of a consonant with \textbf{\mr{ि}}. You have the independent \textbf{\mr{अ}} and the independent \textbf{\mr{उ}}. The independent short \emph{i} is still missing, and a word that starts with it cannot be written without it.
\end{quote}

\subsection*{Script}
\begin{quote}
\mr{इ}
\end{quote}

\textbf{\mr{इ}} is the letter form of the same short \emph{i}. Devanagari's rule has not changed since chapter five: a \textbf{mark} when a consonant carries the vowel, a \textbf{letter} when nothing does. \textbf{\mr{अ}}, \textbf{\mr{आ}}, \textbf{\mr{उ}}, \textbf{\mr{ऊ}}, \textbf{\mr{ए}} are already yours; this is the sixth.

You need it for exactly one reason right now, and it is a good one. The word you learned last chapter points AWAY from you. Its near partner begins with this vowel, and until now the book could say \textbf{\mr{तिथे}} and not its opposite.

\subsection*{Writing — guided copy}
Copy \textbf{\mr{इ}} twice with the model visible. Then write the short-\emph{i} mark beside it and say which of the two can start a word.

\begin{tcolorbox}[breakable,colback=teal!4,colframe=teal!35!black,title={Before you move on}]
Which form opens a word? (\textbf{The letter}.)

\begin{itemize}
\raggedright
  \item \emph{Recall:} say \textbf{\mr{कारण}} and give one reason for anything at all
  \item \emph{Recall:} say the line of your message that means \emph{it was a pleasure to meet you}
\end{itemize}

Source: \href{https://commons.wikimedia.org/wiki/File:Devanagari\_\%E0\%A4\%87\_stroke\_order.svg}{Devanagari i stroke order}.
\end{tcolorbox}
\section[ithe]{\mr{इथे} (ithe) — here}

\label{lesson:MR-C40-ithe}



\begin{quote}
The letter from the last lesson has one job waiting for it.
\end{quote}

\subsection*{You'll want to know: \mr{इथे}}
\begin{quote}
\textbf{\mr{इथे}} — \emph{ithe} — \textbf{here}
\end{quote}

\begin{quote}
\textbf{\mr{मी} \mr{इथे} \mr{आहे}.} — \emph{mī ithe āhe.} — \textbf{I am here.}
\end{quote}

Marathi points two ways and only two. English has \emph{here}, \emph{there} and the older \emph{yonder}; Spanish has three, \emph{aquí}, \emph{ahí} and \emph{allí}. Marathi splits the world in half at your own body: \textbf{\mr{इथे}} is where you are, \textbf{\mr{तिथे}} is everywhere else. There is no middle term to learn.

Now put the three columns side by side, because this is the pattern that will keep paying out for the rest of the book:

\begin{quote}
\textbf{\mr{कुठे}} — asks. \textbf{\mr{इथे}} — near. \textbf{\mr{तिथे}} — far.
\end{quote}

One shared ending, three stems: \textbf{\mr{क}-}, \textbf{\mr{इ}-}, \textbf{\mr{त}-}. When you later meet an unfamiliar Marathi word beginning with one of those three, you already know which of the three jobs it is doing.

\subsection*{Your turn}
\begin{itemize}
\raggedright
  \item \emph{Say it:} \emph{ithe} — here
  \item \emph{Say it:} \emph{mī ithe āhe}
  \item \emph{Say it:} \emph{kuṭhe, ithe, tithe} — with a comma's pause between each
  \item \emph{Recall:} say the line of your message that means \emph{I know Marathi}
\end{itemize}

\begin{tcolorbox}[breakable,colback=teal!4,colframe=teal!35!black,title={Before you move on}]
How many pointing words does Marathi have for place, and where is the line drawn? (\textbf{Two}; \textbf{at your own body}.)

Sources: \href{https://en.wiktionary.org/wiki/\%E0\%A4\%87\%E0\%A4\%A5\%E0\%A5\%87}{Wiktionary: \mr{इथे}}.
\end{tcolorbox}
\section[koṇ · kitī · kuṭhe · kadhī]{Ask your own message six questions}

\label{lesson:MR-R40-asking-column}



\begin{quote}
Six asking words are now yours: \textbf{\mr{काय}}, \textbf{\mr{कसा}}, \textbf{\mr{कोण}}, \textbf{\mr{किती}}, \textbf{\mr{कुठे}}, \textbf{\mr{कधी}}. Name them, then name the one letter they all begin with.
\end{quote}

\subsection*{Your turn}
You wrote a message about yourself many chapters ago and could not question a line of it. Now you can. Take each line and turn it into the question it answers, changing one word and moving nothing.

\begin{itemize}
\raggedright
  \item \emph{Say it:} \emph{tumchaṁ nāv kāy āhe?} — and answer it
  \item \emph{Say it:} \emph{tumhī kase āhāt?} — and answer it
  \item \emph{Say it:} \emph{tumhī kuṭhe rāhtā?} — and answer it
  \item \emph{Recall:} say the line of your message that means \emph{I read Marathi}
  \item \emph{Recall:} write \textbf{\mr{इथे}}, then \textbf{\mr{तिथे}}, and check which one needed a whole vowel letter
\end{itemize}

\begin{tcolorbox}[breakable,colback=teal!4,colframe=teal!35!black,title={Before you move on}]
You can now ask for a person, a number, a place and a time, and answer a place question with \textbf{\mr{इथे}} or \textbf{\mr{तिथे}} without naming anywhere at all. Which two signs did this chapter and the last one give your hand? (\textbf{\mr{थ}} and \textbf{\mr{इ}}.)
\end{tcolorbox}
